\section{Measure a delta before selecting its record}
The previous image now has a defined meaning: it is the source of the last
successfully encoded record. A new frame can be compared with that image using
the XOR kernel developed at the beginning of this chapter. Zero XOR words identify
unchanged pixels. Nonzero words contain the RGB differences to serialize. The
changed-word count is useful evidence, but it does not include command headers or
predict the encoded byte count by itself.

Implement \texttt{video\_encode\_record} with the same four arguments as the full
encoder: source, exact source byte count, destination and native context. The
context grows to 96 bytes. The first 64 bytes retain their earlier field meanings,
except that kind and encoded count describe the selected full or delta record.
At entry the destination capacity must be at least twenty bytes, the smallest
possible record. The actual selected record must fit before any output store.

\begin{table}[htbp]
\centering
\caption{Additional fields in the record encoder context.}
\label{tab:video-delta-encode-context}
\begin{tabular}{r p{1.5in} p{2.8in}}
\toprule
Offset & Field & Meaning \\
\midrule
64 & XOR scratch address & Four-byte-aligned pointer to 4096 native words. \\
72 & XOR scratch capacity & At least 16\,384 bytes. \\
80--95 & Retained suffix & Leave these sixteen bytes unchanged. \\
\bottomrule
\end{tabular}
\end{table}

The caller supplies disjoint source, previous image, scratch, byte destination,
context and private stack windows. Source and previous-image capacities and native
alignment follow the full encoder. Admit scratch alignment and capacity on every
call, including required full frames. Retain the saved registers, frame pointer and
stack pointer, and maintain sixteen-byte stack alignment across nested calls.
Source words must have zero high bytes. For a nonperiodic frame, validate the old
previous-image words too, even if a full representation is eventually selected.
At a required full-frame boundary, old image contents are irrelevant and need not
be read. This allows frame zero to establish fresh prediction state.

A sequence divisible by the positive period requires a full record. Other indices
may use a full or a valid delta record. Completed grading compares record structure,
CRC, sequence and reconstructed numeric colors; it accepts additional full frames
and equivalent delta command boundaries, including valid zero XOR triples. Capacity
12\,304 always suffices for a valid input because a full representation fits.
A smaller buffer may succeed with a compact choice or fail unchanged if the chosen
representation does not fit. A capacity failure does not assert that no possible
encoding could fit. State a course's compression-specific requirements separately
from this general representation contract.

The guided implementation chooses a reproducible policy. After validating colors,
it computes XOR words into scratch. It scans maximal zero runs and maximal nonzero
runs, charging four bytes per command and three per changed pixel. It chooses a
full record when that predicted payload is at least 12\,288 bytes, including an
exact tie. Otherwise it admits sixteen header bytes plus the predicted payload,
then repeats the same run walk to emit it. The two walks must agree exactly.
Scratch is temporary and may change on failure; output, previous image and final
context still must receive no stores on completed failure one.

\begin{table}[htbp]
\centering
\caption{Greedy choices for the supplied frame pairs at nonperiodic sequence one.}
\begin{tabular}{l r r r l}
\toprule
Pair & Changed pixels & Commands & Delta bytes & Chosen record bytes \\
\midrule
Few & 4 & 3 & 24 & 40, delta \\
Still & 0 & 1 & 4 & 20, delta \\
Tail & 1 & 1 & 7 & 23, delta \\
All change & 4096 & 1 & 12\,292 & 12\,304, full \\
Alternating & 2048 & 2049 & 14\,340 & 12\,304, full \\
Tie & 4092 & 3 & 12\,288 & 12\,304, full \\
\bottomrule
\end{tabular}
\end{table}

The few-change pair changes pixel zero by \texttt{0x000001}, pixels 325 and 326
by \texttt{0x102040}, and pixel 4095 by \texttt{0x008000}. Its commands are
skip zero/change one, skip 324/change two, and skip 3768/change one. Those counts
cover all 4096 pixels without a terminator. The payload stores XOR triples, not
current-frame colors. At sequence one, its complete header is
\texttt{01 00 00 00 00 00 00 01 00 00 00 18 D7 F7 3B 51}.
An unchanged frame instead has the four-byte payload \texttt{10 00 00 00}:
skip 4096, change zero. A zero/zero pair would make no progress and is invalid.

Attach this stage with \texttt{attach.py --source encoder-work.asm --kernel-source
xor-work.asm --output record-encoder.asm} from
\texttt{drafts/video-delta-encode-construction}. Supply your earlier XOR working
file as well as the full encoder. The guided importer retains its data and example,
namespaces local labels with \texttt{deltaKernel\_}, and keeps the public XOR entry
names. Instruction names, registers, strings and comments retain their contents;
local symbolic references follow their renamed labels. Use ordinary identifier
labels in this guided import. Free-writing projects may arrange their libraries
differently while preserving the callable contract.

A provided \texttt{video\_record\_crc} leaf helper applies the earlier CRC recurrence
to twelve header bytes followed by a variable-size payload, excluding the checksum
field. It assumes those windows have already been admitted, writes the four CRC
bytes, and returns status zero with the numeric checksum in \texttt{\$v1}. The
new work is choosing and emitting a record; the helper avoids rederiving the same
CRC loop. The supplied caller copies the old previous image to a separate decoder
image before encoding, then decodes the produced bytes against that actual old
state. Copying the newly committed previous image into the decoder would conceal
a broken prediction chain.
